%% ════════════════════════════════════════════════════════════════════════════
\section{Implementation}
\label{sec:implementation}
%% ════════════════════════════════════════════════════════════════════════════

\subsection{Development Methodology and Task Management}
\label{subsec:methodology}

The project followed an Agile-inspired iterative delivery model,
organised into 22 numbered phases spanning hardware assembly, firmware,
backend services, frontend, Docker deployment, and AI analytics.  Each
phase had a clear definition of done and was tracked on a shared Notion
board.  Figure~\ref{fig:notion} shows the task board at the time of
final submission.

\begin{figure}[H]
  \centering
  \includegraphics[width=\linewidth]{screenshot_notion}
  \caption{Notion project board: phase breakdown, assignee tracking,
    and completion status.}
  \label{fig:notion}
\end{figure}

Table~\ref{tab:phases} shows the 22 delivery phases and their
completion status.

\begin{table}[H]
\centering
\caption{Phase Completion Status}
\label{tab:phases}
\begin{tabular}{@{}p{0.6cm}p{6.6cm}p{0.5cm}@{}}
\toprule
\textbf{Ph.} & \textbf{Description} & \textbf{Done} \\
\midrule
0  & Project scaffolding, Docker Compose, env setup         & \checkmark \\
1  & ESP32 firmware — sensors + MQTT publish                & \checkmark \\
2  & Backend foundation — FastAPI, DB models, MQTT bridge   & \checkmark \\
3  & Face recognition service + enrollment API              & \checkmark \\
4  & Session management + attendance engine                 & \checkmark \\
5  & Control API + alert engine                             & \checkmark \\
6  & Moodle sync service                                    & \checkmark \\
7  & WebSocket live streaming                               & \checkmark \\
8  & React frontend                                         & \checkmark \\
9  & Integration testing + documentation                    & \checkmark \\
10 & Full Docker deployment — nginx, Mosquitto wiring       & \checkmark \\
11 & Mock data fallback — backend + frontend demo mode      & \checkmark \\
12 & Raspberry Pi setup runbook                             & \checkmark \\
13 & Demo data seed script                                  & \checkmark \\
14 & JWT auth, professors table, role-based filtering       & \checkmark \\
15 & Docker image slim — stub mode, self-migrating seed     & \checkmark \\
16 & Professor Dashboard — sensor sparklines, session list  & \checkmark \\
17 & Dashboard bug fixes — enrolled count, key resets       & \checkmark \\
18 & Full frontend visual redesign — CSS token system       & \checkmark \\
19 & At-Risk Explanation — Ollama/phi3-mini pipeline        & \checkmark \\
20 & At-Risk pipeline performance optimisation              & \checkmark \\
21 & LLM-Assisted Forecasting + Recharts trend chart        & \checkmark \\
22 & Snapshot-per-cycle attendance; laptop mode conflict fix& \checkmark \\
\bottomrule
\end{tabular}
\end{table}

\subsection{Hardware Assembly}
\label{subsec:hardware}

Table~\ref{tab:hardware_bom} lists the bill of materials.  All
components interface to either the ESP32 or the Raspberry Pi~4B.

\begin{table}[H]
\centering
\caption{Hardware Bill of Materials}
\label{tab:hardware_bom}
\begin{tabular}{@{}p{2.5cm}p{4.8cm}@{}}
\toprule
\textbf{Component} & \textbf{Role} \\
\midrule
Raspberry Pi~4B 4~GB  & Central hub; runs all Docker services and camera \\
RPi Camera Module~v2  & 8~MP camera input for face recognition \\
ESP32 WROOM-32        & Sensor node + relay controller via MQTT \\
DHT21 (AM2301)        & Temperature and humidity (GPIO~4) \\
MQ-135                & Air quality / CO$_2$ proxy (ADC GPIO~34) \\
ACP014 Sound Sensor   & Binary occupancy/noise detection (GPIO~35) \\
4-Channel Relay       & AC (CH1), lighting (CH2), spare (CH3/4) \\
LCD 16$\times$2 I2C   & Local room status display (I2C GPIO~21/22) \\
Logic Level Converter & 3.3\,V $\leftrightarrow$ 5\,V: ESP32 to relay \\
USB-C 5\,V 3\,A PSU  & Raspberry Pi power supply \\
32~GB MicroSD         & RPi OS and Docker volume storage \\
Breadboard + wires    & Prototyping connections \\
\bottomrule
\end{tabular}
\end{table}

The critical wiring consideration is the logic-level converter between
the ESP32 (3.3\,V logic) and the relay module (5\,V TTL inputs).
Driving the relay directly from the ESP32 GPIO would result in
unreliable switching, as the relay's opto-isolator threshold is
nominally 5\,V.  The converter's low-voltage side connects to ESP32
GPIO~26--27 (relay CH1/CH2), and the high-voltage side feeds the relay
IN pins.

Figure~\ref{fig:rpi4b} shows the Raspberry Pi~4B with the camera
module attached.

\begin{figure}[H]
  \centering
  \includegraphics[width=0.85\linewidth]{photo_rpi4b}
  \caption{Raspberry Pi~4B with Camera Module~v2 mounted.  The RPi
    hosts all backend services and the face recognition loop.}
  \label{fig:rpi4b}
\end{figure}

\subsection{Firmware Development (ESP32)}
\label{subsec:firmware}

The Arduino sketch (\texttt{classroom\_node.ino}) implements a
polling loop: every 5~seconds it reads all sensors, constructs JSON
payloads, and publishes each to its MQTT topic.  Configuration
parameters (WiFi SSID/password, MQTT broker IP, room ID, threshold
values) are centralised in \texttt{config.h} so deployment to a new
classroom requires only editing that file and re-flashing.

MQTT reconnection is handled by a \texttt{reconnect()} function called
from \texttt{loop()} whenever \texttt{client.connected()} returns false,
with a 5-second back-off delay between attempts.  The LCD displays live
temperature, humidity, air quality, and sound status on alternating
lines, providing local room status without requiring the professor to
access the dashboard.

\subsection{Backend Development}
\label{subsec:backend_dev}

\textbf{Project structure.}  The backend follows a flat service
architecture under \texttt{backend/app/}: \texttt{api/} for route
handlers, \texttt{services/} for business logic, \texttt{models/} for
SQLAlchemy ORM models and Pydantic schemas, and top-level
\texttt{main.py}, \texttt{config.py}, \texttt{database.py}, and
\texttt{redis\_client.py}.

\textbf{Database initialisation.}  Alembic auto-migration runs on
startup via a subprocess call (\texttt{alembic upgrade head}) inside the
FastAPI lifespan context manager.  This ensures the schema is always
current without a separate migration step.

\textbf{MQTT bridge.}  The \texttt{aiomqtt} bridge runs as an
\texttt{asyncio.Task}, maintaining a persistent subscription to
\texttt{classroom/+/sensors/\#} and \texttt{classroom/+/status}.
Exponential back-off reconnection (initial 1~s, max 60~s) handles
Mosquitto restarts.

\textbf{Environment variables.}  Table~\ref{tab:env_vars} lists the
primary configuration variables.

\begin{table}[H]
\centering
\caption{Primary Environment Variables}
\label{tab:env_vars}
\begin{tabular}{@{}p{3.2cm}p{4.3cm}@{}}
\toprule
\textbf{Variable} & \textbf{Description and Default} \\
\midrule
\texttt{DATABASE\_URL}            & PostgreSQL async connection string \\
\texttt{REDIS\_URL}               & Redis connection URL \\
\texttt{MQTT\_BROKER\_HOST}       & Mosquitto hostname (default: \texttt{mosquitto}) \\
\texttt{ROOM\_ID}                 & Active room identifier (default: \texttt{room1}) \\
\texttt{MOCK\_MODE}               & \texttt{true}: synthetic MQTT without ESP32 \\
\texttt{FACE\_RECOGNITION\_ENABLED} & \texttt{true}: real camera + DeepFace (RPi only) \\
\texttt{SECRET\_KEY}              & JWT signing key — must be replaced in production \\
\texttt{REQUIRE\_AUTH}            & \texttt{false} in dev; \texttt{true} in production \\
\texttt{AT\_RISK\_THRESHOLD}      & Attendance rate below which a student is at-risk (default: 0.70) \\
\texttt{ATTENDANCE\_CYCLE\_DURATION} & Seconds per scan-evaluate cycle (default: 60) \\
\texttt{OLLAMA\_BASE\_URL}        & Ollama service URL (default: \texttt{http://ollama:11434}) \\
\texttt{FORECAST\_WINDOW}         & Past sessions included in trend (default: 8) \\
\bottomrule
\end{tabular}
\end{table}

\subsection{Face Recognition Pipeline}
\label{subsec:recognition}

Figure~\ref{fig:face_recog} shows the complete enrollment and
recognition flow.

\begin{figure}[H]
  \centering
  \includegraphics[width=\linewidth]{fig10_face_recognition_flow}
  \caption{Enrollment and recognition pipeline: from image upload to
    attendance record creation.}
  \label{fig:face_recog}
\end{figure}

\textbf{Enrollment.}  The professor uploads up to five reference images
per student via the Enrollment page (Figure~\ref{fig:enrollment}).
Each image is passed to DeepFace's \texttt{represent()} function with
the FaceNet backend, producing a 128-dimensional float32 embedding.
The per-image embeddings are averaged and stored as a BYTEA column in
the \texttt{face\_encodings} table.  No raw images are retained.

\begin{figure}[H]
  \centering
  \includegraphics[width=\linewidth]{screenshot_enrollment}
  \caption{Enrollment page: student registration and face image
    upload.}
  \label{fig:enrollment}
\end{figure}

\textbf{Recognition loop — snapshot-per-cycle model.}  While a session
is active, a background \texttt{asyncio.Task} runs the recognition loop:
\begin{enumerate}
  \item \textbf{SCAN phase} (duration: \texttt{ATTENDANCE\_CYCLE\_DURATION}
    seconds at \texttt{RECOGNITION\_FPS} fps): capture frames, compare
    each detected face embedding against all stored embeddings using
    cosine distance.  Any student whose distance falls below 0.40 is
    added to \texttt{seen\_this\_cycle}.

  \item \textbf{EVALUATE phase} (single transaction): for every enrolled
    student --- if in \texttt{seen\_this\_cycle}: mark present; if not:
    mark absent.  Records where \texttt{adjusted\_by IS NOT NULL} are
    never touched.
\end{enumerate}

This bidirectional model captures students who leave mid-session, which
the original one-shot model missed.

\subsection{Frontend Development}
\label{subsec:frontend_dev}

The frontend is scaffolded with Vite and React~18.  The
\texttt{useLiveSensors} hook manages the WebSocket lifecycle, the demo-mode
watchdog, and message dispatch to their respective React state slices.
If no WebSocket data arrives within 8~seconds of connecting, the hook
activates a client-side mock generator and renders the
\texttt{DemoModeBanner}, allowing professors to review the interface
without a running backend.

Recharts \texttt{AreaChart} with \texttt{linearGradient} fill renders
the sensor sparklines on the dashboard.  The forecasting page adds a
70\% dashed \texttt{ReferenceLine} and optionally appends a projected
``Next'' data point to the trend series.

\subsection{Infrastructure and Deployment}
\label{subsec:infra}

The full stack is containerised with Docker Compose.  The backend
container builds from a custom image with DeepFace and TensorFlow
excluded (stub mode only), reducing the image from approximately 3~GB
to under 800~MB.  Named Docker volumes persist PostgreSQL data, Redis
data, Mosquitto broker state, and the Ollama model cache across
container restarts.

The nginx container serves the React static files on port~3000 and
proxies all \texttt{/api/*} and \texttt{/ws/*} requests to the FastAPI
backend at port~8000, making the backend never directly exposed to the
LAN.

\subsection{Seed Data and Demo Environment}
\label{subsec:seed}

The \texttt{seed.py} script runs Alembic migrations first, then
populates the database with 35 students, 6 courses, 5 professors, and
30 past sessions with synthetic attendance records.  Running the stack
with \texttt{MOCK\_MODE=true} activates the sine-wave sensor publisher,
making all dashboard charts and sparklines live without an ESP32.
Figure~\ref{fig:swagger} shows the FastAPI Swagger UI used during
backend development and integration testing.

\begin{figure}[H]
  \centering
  \includegraphics[width=\linewidth]{screenshot_swagger}
  \caption{FastAPI Swagger UI (\texttt{/docs}): interactive API
    explorer used for endpoint validation during development.}
  \label{fig:swagger}
\end{figure}
